% 第 4 题程序框图（累减计数型，双菱形双汇入结构）：按固定源 PDF 逐框重绘（10 框；1 pt → 0.04 cm）。
% 结构要点（易错，矢量路径逐段核对）：
%   主干 x=0（两个菱形与主线上全部处理框同一中心线）；左支 x=-2.45（S=2S）；右回线 x=+2.34。
%   第一菱形「否」自左顶点向左折下入 S=2S，其回接自框底垂直下、水平右，汇入主干 y=4.87
%   （S=S-6 之后、n=n+1 之前）；第二菱形「否」自右顶点向右、向上至 y=7.45、向左汇入主干
%   （n=1 之后、第一菱形之前）。两处汇入高度不同，均不做自环。
%   全图 11 个箭头；输出框为向右倾平行四边形；无输入框。
\begin{tikzpicture}[>=stealth,every node/.style={font=\small},
  term/.style={draw,rounded corners,minimum width=1.5cm,minimum height=5.5mm},
  box/.style={draw,minimum height=5.5mm,inner xsep=4pt},
  io/.style={draw,trapezium,trapezium left angle=70,trapezium right angle=110,minimum width=1.8cm,minimum height=5.5mm},
  decision/.style={draw,diamond,aspect=4,minimum width=3.2cm,inner sep=0pt}]
\node[term] (start) at (0,10.36) {开始};
\node[box] (s4) at (0,9.16) {\(S=4\)};
\node[box] (n1) at (0,8.02) {\(n=1\)};
\node[decision] (t1) at (0,6.75) {\(S\geqslant 6?\)};
\node[box] (s6) at (0,5.46) {\(S=S-6\)};
\node[box] (s2s) at (-2.45,5.46) {\(S=2S\)};
\node[box] (np1) at (0,4.28) {\(n=n+1\)};
\node[decision] (t2) at (0,3.03) {\(n>3?\)};
\node[io] (out) at (0,1.74) {输出 \(S\)};
\node[term] (end) at (0,0.46) {结束};
\draw[->] (start)--(s4);
\draw[->] (s4)--(n1);
\draw[->] (n1)--(t1);
\draw[->] (t1.south)--(s6.north);
\draw[->] (s6)--(np1);
\draw[->] (np1)--(t2);
\draw[->] (t2.south)--(out.north);
\draw[->] (out)--(end);
% 第一菱形「否」：左顶点向左折下入 S=2S
\draw[->] (t1.west) -| (s2s.north);
% S=2S 回接：框底垂直向下、水平向右汇入主干（汇入点高于 n=n+1）
\draw[->] (s2s.south) |- (0,4.87);
% 第二菱形「否」：右顶点向右、向上、向左汇入主干（汇入点高于第一菱形）
\draw[->] (t2.east) -- (2.34,3.03) -- (2.34,7.45) -- (0,7.45);
% 分支标注（「否」在各自水平线上方、「是」在主干竖线右侧）
\node at (-1.64,7.05) {否};
\node at (0.65,6.15) {是};
\node at (1.55,3.38) {否};
\node at (0.65,2.42) {是};
\end{tikzpicture}
